We study residue-class transitions of consecutive primes modulo 30 as a finite computational dynamical system. Using the first-order transition operator P as a Markov baseline, we compare the empirical two-step operator P^{(2)} against P^2 and analyze the residual operator Delta = P^{(2)} - P^2. Across the computed range, the residual exhibits low-rank structure, stable singular modes, and measurable prediction improvements under low-rank correction. Synthetic controls, including iid shuffles, Markov-generated sequences, block shuffles, balanced shuffles, and gap-based shuffles, show that the observed structure is not explained by marginal residue frequencies or first-order transition statistics alone. Bootstrap and null validation indicate statistically meaningful separation from structure-destroying controls while remaining consistent with local block-preserving resampling. These results provide empirical computational evidence for higher-order residue memory in finite prime-gap transition data, without claiming an asymptotic theorem.
